\subsection{The dihedral genus-two quotients: the complete fixed subgroup}

\begin{statement}{A.39}{Proved}{Every dihedral genus-two first quotient has exactly cusp fixed subgroups}
Let
\[
 (a,b)\in\{(0,-1),(6,1),(-6,1)\},\qquad c^4=b,
\]
and let
\[
 D_{a,c}:z^2=(x^2-4c)(x^4-4cx^2+a+2c^2),
 \qquad Y_{a,c}=D_{a,c}\setminus\{Q_+,Q_-\},
\]
where \(Q_+,Q_-\) are its two points at infinity.
These include both reflection conjugacy classes for each
of the three \(D_8\) covers in A.4. For every prime
\(\ell\), every original packet of Part A, and every
permitted open subgroup \(\Delta_0\) preserving the
curve and its selected cusp lift \(P\),
\[
 \left(\pi_1(Y_{a,c,\bar K})^{(\ell)}\right)^{s_b(\Delta_0)}
                      =I_{P,\ell}.                    \tag{A.39.1}
\]
At odd primes the same argument works for every cusp
packet over a finite number field of definition, with
that field's own derived Galois subgroup.

The proof uses the actual pointed reduction at two,
an unbounded weighted monodromy kernel, and the
cyclotomic criterion of A.38. It does not transfer an
embedding through the unramified double cover
\(C_4\to D_{a,c}\) at an odd prime, nor require a
primitive Collino regulator for the conjugate cusps.
\end{statement}

\begin{proof}
\paragraph{The two cusps meet in the unpointed stable model.}
The proper Jacobian is isogenous to \(E\times F\),
where the complementary \(j\)-invariant is \(8000\),
\(10976\), or \(128\), by A.4. Together with
\(j(E)=1728\), all have positive valuation at two.
After a finite extension the Jacobian has good
supersingular reduction. The unpointed stable fibre
is of compact type and has \(2\)-rank zero.
A stable unmarked compact-type genus-two curve is
either a smooth genus-two curve or two smooth
elliptic curves joined at a node. It has no rational
components: a rational vertex in a stable tree needs
at least three branches, each eventually containing
positive genus, which would give total genus at
least three.

Write \(\delta=[Q_+-Q_-]\). A.4 proves
\(4\delta=0\), in fact exact order four, for every
\(a,c\) under consideration. Only the displayed
annihilation is needed here. A useful uniform
identity behind it is
\[
 \begin{aligned}
 A&=x^3-4cx,\qquad d=a+2c^2,\qquad
 z^2-A^2=d(x^2-4c),\\
 f&=\frac{(z-A)^2}{x^2-4c}=d\frac{z-A}{z+A},\qquad
 \operatorname{div}(f)=4(Q_+-Q_-).
 \end{aligned}                                      \tag{A.39.2}
\]
Here \(d\ne0\) in all the indicated cases.
Since the geometric special Jacobian has \(2\)-rank
zero, it has no nonzero geometric points of order
a power of two. The specialization of \(\delta\)
is therefore zero.

Take the Abel map defined by either cusp as its
base point. By
\cite[Theorem 1.1]{fghDupuyRabinoff}, it extends to
the compact-type model and is a closed immersion
because the unpointed special fibre has no rational
components. Thus \(Q_+\) and \(Q_-\) specialize to
the same point of that fibre. For the two-component
case this can also be read directly from
\cite[Proposition 3.2]{fghDupuyRabinoff}: the special
Abel map embeds the two elliptic components as the
two Jacobian axes, meeting only at the origin.
In particular the assertion includes specialization
to a node, not just to the smooth locus.

If the unpointed fibre is smooth, separating the two
sections attaches a rational component carrying both
cusps and a single node. The edge to the genus-two
component has strictly positive length. Contracting
unmarked rational subdivisions merely adds their
lengths.

If the unpointed fibre is \(E_1\cup E_2\), consider
the hyperelliptic involution \(h(x,z)=(x,-z)\).
It exchanges the two cusps and acts as \(-1\) on
proper homology. It cannot exchange the elliptic
components, since an involution exchanging two
two-dimensional homology spaces has a two-dimensional
invariant space. Thus it preserves both components
and acts as their elliptic involution with origin
the common node. Each supersingular elliptic
component has only this one geometric fixed point.
The common specialization of \(Q_+,Q_-\), being
fixed by \(h\), must consequently be the node.

The involution preserves the two branches of that
node, hence fixes the metric edge pointwise. Its
exchange of the two cusp sections shows that their
interior retraction points agree. Subdivide at that
point after a finite extension. On the inserted
rational bridge, \(h\) fixes the two attachments.
In characteristic two a nonidentity involution of
\(\mathbf P^1\) has only one geometric fixed point,
so this bridge action is the identity. The two
sections therefore still meet at the same smooth
point of the bridge. Separating them attaches a
rational two-mark leaf, with no genus in that
branch. In the stable pointed model the other end
of its attaching edge is the rational bridge
joining the two elliptic components. Denote the
leaf edge length by \(r>0\), and the two arm lengths
by \(m_1,m_2>0\). Their positivity follows because
the common retraction is interior to the original
node edge and the cusp sections collide on its
inserted bridge.

\paragraph{Actual weighted monodromy.}
For odd \(\ell\), give the free rational Lie algebra
of the punctured genus-two curve its weight filtration:
the four proper handle generators have weight one,
and the logarithm \(t\) of the selected cusp has
weight two. Write
\[
 \omega_1=[a_1,b_1],\quad\omega_2=[a_2,b_2],
 \qquad\Omega=\omega_1+\omega_2.
\]
Use an ordinary base point reducing to the rational
two-mark leaf. Its other cusp and the node form the
other two boundaries of that rational tripod.
The node encloses all proper genus and has first
weight two boundary class \(\Omega\), up to the
common orientation convention. The based affine
edge-twist formula in
\cite[Definition 3.1.5 and Theorem 3.1.6]{fghBettsDogra}
therefore gives the following first nonzero weighted
derivation. In the smooth genus-two case it is
\[
                 Nt=0,\qquad Nv=r[\Omega,v]
                  \quad(v=a_1,b_1,a_2,b_2),\quad r>0.
                                                               \tag{A.39.3}
\]
In the two-elliptic case it is
\[
             Nt=0,\qquad
             Nv_i=[r\Omega+m_i\omega_i,v_i]
                  \quad(v_i=a_i,b_i),\quad r,m_i>0.
                                                               \tag{A.39.4}
\]
The common leaf edge contributes \(r\Omega\) to
both handle blocks, while each arm contributes its
own \(m_i\omega_i\). All these terms have weight
shift two. There is no weight-one edge term.

The ordinary-point comparison with the chosen
tangent takes place through the selected cusp's
punctured formal residue disc, exactly as in A.36
and A.38. Moving the ordinary point out along the
cusp half-edge contributes only a cusp inner term.
The comparison cocycle and every completed tangent
packet change lie in the actual selected inertia.
Thus the genuine derivation is
\[
                         D_\lambda=N+\lambda\operatorname{ad}t,
                   \qquad\lambda\in\mathbf Q_\ell.   \tag{A.39.5}
\]
Equivalently, a weight-two inner discrepancy fixing
\(t\) has conjugator in its centralizer
\(\mathbf Q_\ell t\); no weight-one discrepancy
can fix \(t\).

Choose a finite global field \(k\) defining the data
and realizing the indicated local semistable model,
with \(\Gamma=\overline{[G_k,G_k]}\subset\Delta_0\).
An actual local Frobenius--tame-inertia commutator
at two belongs to \(\Gamma\), and its first weighted
derivation is a nonzero common multiple of (A.39.5).
Indeed its tame parameter is multiplied by
\(q_{\rm res}-1\ne0\). The common factor can be
divided out in the rational kernel calculation.
No nonvanishing condition on \(\lambda\) is assumed.

\paragraph{The smooth-component all-degree kernel.}
First use (A.39.3). The leading associative word of
\(\Omega\), in a degree-lexicographic order with
\(a_1>b_1>a_2>b_2\), is \(a_1b_1\), with coefficient
one and no proper self-overlap. The triangular
bimodule construction of A.27 therefore applies to
\(A_V=\mathbf Q_\ell\langle a_1,b_1,a_2,b_2\rangle\)
and \(\Omega\). Specifically, let \(\mathcal C\)
be all nonempty handle words which neither begin
nor end with \(a_1b_1\). Stripping initial and
terminal powers of that word gives the basis
\[
 \{\Omega^n:n\ge0\}\ \sqcup\
 \{\Omega^u w\Omega^v:w\in\mathcal C,\ u,v\ge0\}.
                                                               \tag{A.39.6}
\]
Each element has its distinct indicated leading
word, so this is a basis in every finite homogeneous
degree. Adjoin the free letter \(t\), and set
\(B=\mathbf Q_\ell\langle t,\Omega\rangle\).
A word outside \(B\) has a first nonboundary handle
core in (A.39.6), with an arbitrary \(B\)-prefix.
On that prefix the part of \(D_\lambda\) increasing
its length is
\[
                          \lambda L_t+rR_\Omega.
                                                               \tag{A.39.7}
\]
It is injective: if \(\lambda=0\), use right
multiplication by \(\Omega\); otherwise count
occurrences of \(t\). Subsequent terms cannot
remove the first core and preserve its prefix
length. The kernel is therefore contained in
\(B\). There the derivation is
\(\lambda\operatorname{ad}t\), so its associative
kernel is \(B\) for \(\lambda=0\) and
\(\mathbf Q_\ell[t]\) otherwise.

\paragraph{The two-component all-degree kernel.}
For (A.39.4), use the two separate handle-core
bases of A.27 and put
\(B=\mathbf Q_\ell\langle t,\omega_1,\omega_2\rangle\).
On this freely generated algebra,
\[
 N_Bt=0,\qquad N_B\omega_i=r[\Omega,\omega_i],
 \qquad D_\lambda|_B=N_B+\lambda\operatorname{ad}t.
                                                               \tag{A.39.8}
\]
For a first core of type \(i\), the part increasing
the length of its \(B\)-prefix is the operator
\[
 K_i=N_B+\lambda L_t+R_{r\Omega+m_i\omega_i}.
                                                               \tag{A.39.9}
\]
This follows by combining \(D_\lambda\) on the
prefix with the left half of
\([r\Omega+m_i\omega_i+\lambda t,w]\) on the
core; the right \(\lambda t\) terms cancel.
All other terms preserve the prefix length and
the first core.

If \(\lambda\ne0\), \(N_B\) and both right
multiplications preserve the number of \(t\)'s,
while \(\lambda L_t\) raises it. The largest
\(t\)-count proves injectivity of \(K_i\).
If \(\lambda=0\), use the free generating system
\(\Omega,\omega_i,t\) of \(B\).
In this system \(N_B\Omega=N_Bt=0\) and
\(N_B\omega_i=r[\Omega,\omega_i]\), so \(N_B\)
preserves \(\omega_i\)-count. Right multiplication
by \(r\Omega\) also preserves it, whereas
\(m_iR_{\omega_i}\) raises it. Since \(m_i\ne0\),
\(K_i\) is again injective. Taking a greatest
prefix length proves that no element outside
\(B\) belongs to the kernel.

It remains to compute the kernel in \(B\), where
\(\Omega,t,\omega_1\) are free generators.
On \(B_0=\mathbf Q_\ell\langle t,\Omega\rangle\)
the derivation is \(\lambda\operatorname{ad}t\),
and
\[
                 D_\lambda\omega_1
                           =[r\Omega+\lambda t,\omega_1].
\]
The first-slot argument of A.36 now has operator
\(\lambda L_t+rR_\Omega\), which is injective.
Thus the associative kernel is \(B_0\) if
\(\lambda=0\) and \(\mathbf Q_\ell[t]\) otherwise.

In both geometric cases, all displayed boundary
generators are primitive in the tensor Hopf
algebra. Intersecting their associative kernels
with primitives gives the exact free-Lie formula
\[
 \ker D_\lambda=
 \begin{cases}
     \mathbb L_{\mathbf Q_\ell}(t,\Omega),&\lambda=0,\\
     \mathbf Q_\ell t,&\lambda\ne0.
 \end{cases}                                           \tag{A.39.10}
\]
This holds in all weights; both \(t\) and
\(\Omega\) have weight two.

\paragraph{The actual fixed subgroup.}
Let \(H\) be the actual \(\Gamma\)-fixed subgroup
of the geometric pro-\(\ell\) group. Its unipotent
closure is pointwise \(\Gamma\)-fixed. Every
leading weighted Lie class in that closure is
therefore killed by the actual derivation (A.39.5),
and hence belongs to (A.39.10). The full \(G_k\)
action on \(t\) is \(\chi_\ell\), by the selected
cuspidal packet; its action on \(\Omega\) is the
same character, by the symplectic similitude on
proper homology. The closure has no odd weight
pieces, and its weight \(2m\) piece has scalar
character \(\chi_\ell^m\).

The criterion of A.38 applies, giving \(H=I_{P,\ell}\).
Since \(\Gamma\subset\Delta_0\) and the original
\(\Delta_0\) fixes the selected inertia, this proves
(A.39.1) for every odd prime and every stated packet.
At \(\ell=2\), A.22 applies to each reflection
subgroup, whose Galois closure is the original
\(D_8\) cover. This gives the remaining prime.

The exact certificate \texttt{code/a\_D4\_fixed.py}
checks (A.39.2), the two-component derivation
identities on \(B\), the triangular basis used for
a smooth genus-two component, and finite weighted
consistency checks for both geometric cases and
several packet scalars. The proof of (A.39.10)
is the unbounded argument above, not an inference
from those finite checks.
\end{proof}
